\section{Method}
\label{sec:method}

\subsection{System Overview}

Security Cuff is a two-layer runtime defense for deployed VLA
policies.
The fast layer performs dense, low-latency screening on rollout
prefixes, and the slow layer performs higher-cost verification only on
the subset that has already been flagged as risky.
This separation is the core design choice of the system.
A single detector is typically forced to trade off recall, latency, and
verification strength within one module, whereas the dual-layer design
assigns broad online coverage to the fast layer and selective
consequence-aware verification to the slow layer.

Given a rollout prefix $\tau_{1:t}$, the fast layer first produces a
stepwise risk score $s_t$.
An online aggregation rule then updates the rollout-level risk estimate
$r_t$ and decides whether the current prefix should be escalated.
If escalation is triggered, the slow layer receives the same rollout
prefix together with the cached runtime context and returns a stronger
verification signal $v_t$.
The final system decision is made by combining the fast screening
signal, the routing state, and the slow verification result.
The remainder of this section describes the fast layer, the slow layer,
and the routing rule in turn.

\subsection{Fast Layer: Rollout-Risk Screening}

The fast layer is designed for high-recall online screening.
Its role is to detect risky rollouts early under tight latency
constraints rather than to recover the full attack mechanism behind
every event.
This requirement is particularly important in our main setting, where
the positive class is backdoor-triggered dangerous execution and
auxiliary out-of-distribution failures are evaluated separately.
In this setting, the fast layer should separate successful execution
from backdoor-induced risk with minimal runtime overhead.

The fast layer operates on runtime-accessible signals exposed by the
deployed policy.
Its core input is the policy hidden state $h_t$, which summarizes the
current observation and execution context.
When available, we also include a short-horizon runtime context
$u_t$, such as recent actions or proprioceptive measurements, to retain
local execution dynamics that may not be fully captured by $h_t$
alone.
To anchor detection to successful execution, we construct a success
reference bank $\mathcal{S}$ from successful training rollouts and
derive a reference summary $\rho_t(\mathcal{S})$ for the current step.
We then form a residual feature
\[
  \delta_t = h_t - \rho_t(\mathcal{S}),
\]
and use the detector input
\[
  \phi_t = [\,h_t;\delta_t;u_t\,].
\]
The fast detector maps $\phi_t$ to a scalar stepwise risk score
\[
  s_t = f_{\mathrm{fast}}(\phi_t) \in [0,1].
\]

This representation is useful because it turns online detection into a
comparison against successful execution rather than a search for a
single abnormal pattern.
The risky class is heterogeneous, but successful rollouts are often
more structurally consistent.
The residual feature therefore provides a stable reference frame for
online screening.
It also avoids requiring full disentanglement of checkpoint, scene, and
behavior factors in the hidden state.
For deployment-time risk detection, any stable runtime signal that
improves separation between successful and risky rollouts is potentially
useful, and the fast layer is designed to exploit that signal under a
lightweight inference budget.

\subsection{Slow Layer: Consequence-Aware Verification}

The slow layer is designed for selective verification after the fast
layer has already raised concern.
Fast screening is necessary for broad coverage, but it does not by
itself provide the strongest possible basis for intervention.
In physical environments, missed detections can lead to irreversible
outcomes, while false alarms can unnecessarily interrupt successful
execution.
The slow layer addresses this tension by spending additional compute
only on the routed subset of suspicious rollouts.

When the system escalates a rollout prefix, the slow layer receives the
current prefix together with the cached runtime context and evaluates
whether continuing execution is likely to produce an unsafe outcome.
We denote its verification output by
\[
  v_t = f_{\mathrm{slow}}(\tau_{1:t}) \in [0,1].
\]
The slow layer may instantiate this verification step with a learned
world model, a short-horizon consequence predictor, or an
attribution-oriented consistency check~\citep{ha2018world,daydreamer}.
The defining property is not the exact verifier family, but the role of
this stage in the system: it provides a higher-cost, higher-fidelity
assessment for a small routed subset rather than dense scoring on every
time step.

This division of labor gives the slow layer a different capability from
the fast layer rather than a slower copy of the same detector.
The fast layer is optimized for recall under strict latency
constraints.
The slow layer is optimized for stronger verification on ambiguous or
high-risk cases.
Together, the two layers support a system-level tradeoff in which most
rollouts are handled cheaply, while difficult cases receive additional
reasoning before the final intervention decision is made.

\subsection{Online Aggregation and Routing}

The fast layer emits stepwise scores, but deployment decisions must be
made at the rollout level.
We therefore maintain an online aggregated risk estimate
\[
  r_t = g(r_{t-1}, s_t),
\]
where $g$ is a lightweight online update rule.
The system continues normal execution while $r_t$ remains below the
escalation threshold $\gamma$.
Once $r_t \ge \gamma$, the current rollout prefix is routed to the slow
layer, which returns the verification score $v_t$.
The final action is then determined by a decision rule
\[
  d_t = H(r_t, v_t),
\]
which maps the current fast-layer risk and slow-layer verification to a
deployment action such as continue, warn, or intervene.

This routing rule couples the two layers into a single system.
It keeps the common-case computation close to the fast detector while
reserving expensive verification for the subset that is most likely to
benefit from it.
As a result, the overall design targets three objectives
simultaneously: early warning on risky rollouts, controlled false
alarms on successful rollouts, and bounded deployment-time overhead.
